% Q10_solutions.tex — REFERENCE SOLUTIONS to paper/Q10_problems.tex.
% Author-only and read-denied in study mode (Read(./paper/solutions/**)).
% Built only by `make -C paper paper-solutions` (scripts/ci.sh, or the
% author's linked worktree); `make paper` never touches this directory.
%
% Every verifier result quoted below was captured on the course pins (Dafny
% 4.11.0, Verus 0.2026.09.27.3cf1832) through `scripts/vcheck --module Q10`,
% on 2026-10-07, one solver job at a time. The four loop files and their
% outputs are archived under
% ~/.cache/dafny-verus-learn-0/author-probes/dv5ak/q10-kit/ (OUTPUTS.txt). B2's
% fields are graded against the learner's own cheatsheets.md; this file
% records which lines of the pinned range each field draws on and never copies
% the card.
\documentclass[11pt]{article}
\input{preamble}

\title{Q10 --- Practice I-A \\ \large Reference solutions (author-only)}
\author{dafny-verus-learn-0}
\date{}

\begin{document}
\maketitle
\sloppy

\section*{How these are graded}
B1 is a concept answer (\texttt{correctness}: the grader runs each of the four
files with \texttt{scripts/vcheck --module Q10} and tries to counterexample
each reason). B2 is a paper card (\texttt{card-fidelity}, per field, and the
ninth line as a short concept answer). The solutions below give the answers;
the rubric never states them.

% ---------------------------------------------------------------------------
\section*{B1 (concept): the four loops}

Captured outputs (\texttt{raw output:} lines dropped):
{\small
\begin{verbatim}
vcheck: a.dfy [dafny, learner loop]: failed — 22 verified, 0 holes, 1 error
a.dfy:4:3: error: a postcondition could not be proved on this return path
a.dfy:3:15: related: this is the postcondition that could not be proved
vcheck: b.dfy [dafny, learner loop]: verified — 23 verified, 0 holes, 0 errors
vcheck: c.dfy [dafny, learner loop]: failed — 22 verified, 0 holes, 1 error
c.dfy:6:5: error: cannot prove termination; try supplying a decreases clause for the loop
vcheck: d.rs [verus, learner loop]: failed — - verified, 0 holes, 1 error, stage resolve
d.rs:10:5: error: loop must have a decreases clause
\end{verbatim}
}

\begin{description}
\item[A.] Fails at verification: the postcondition \dfy{i == |s|}. Termination
  is fine (the guessed measure \dfy{|s| - i} decreases and is bounded while the
  guard holds). The loop has no invariant, so after it the verifier knows only
  the negated guard, \dfy{i >= |s|}, not \dfy{i <= |s|}. The measure is not the
  problem; the missing invariant is.
\item[B.] Verifies. The guessed measure is \dfy{j} (from \dfy{j > 0}); it
  decreases by one each step and is a \dfy{nat}. After the loop,
  \dfy{!(j > 0)} with \dfy{j: nat} gives \dfy{j == 0}.
\item[C.] Fails at verification: termination. The guessed measure for
  \dfy{i != n} is the distance between the two sides; stepping by 2 from 0 can
  jump over an odd \dfy{n}, after which the distance grows and the loop does
  not stop. (The invariant \dfy{i <= n} is not kept for odd \dfy{n} either;
  Dafny reports the termination obligation. Either reason, with the odd case,
  is a full answer.) Here the guessed measure is the loop's problem, and no
  measure fixes it: the loop really does not terminate for odd \dfy{n}.
\item[D.] Stops at \texttt{resolve}: ``loop must have a decreases clause''.
  Verus guesses no measure for an exec loop; the invariant is fine, and adding
  \vrs{decreases v.len() - i} makes it verify (the chapter's \S4 P2 probe has
  that shape).
\item[The measure Dafny uses.] When a loop has no \texttt{decreases}, Dafny
  guesses one from the guard: for \dfy{A < B} (or \dfy{<=}) the difference
  \dfy{B - A}, for \dfy{A > B} the difference \dfy{A - B}, for \dfy{A != B} the
  distance between the sides. A guess is checked like a written measure, so it
  fails exactly when the loop's step does not decrease it toward its bound.
  Accept any statement that the guess is a difference read off the guard.
\end{description}

% ---------------------------------------------------------------------------
\section*{B2 (paper card 02.11): reference fields and their lines}

Pinned range: \texttt{cheatsheets.md} lines 888--913 (computed 2026-10-07),
\[ \texttt{sha256 d9fda575ed350eed76ac951aaca5717f1063acb0f78f146ac4ef1201d74c056a}. \]
A field is correct when it agrees with the lines named; wording is free. Each
entry: the reference content (paraphrased), then the lines it draws on.

\begin{description}
\item[Task shape.] Proof repair of one gap: a real Dafny lemma from which one
  assertion that the proof needs has been deleted; the model must supply an
  assertion (its place and its content) so that the lemma verifies again. The
  benchmark, DafnyGym, has 145 such tasks from three Amazon-maintained
  codebases. \emph{Lines 890, 892.}
\item[Oracle.] The Dafny verifier (4.3.0) accepts the lemma. \emph{Line 892.}
\item[What is trusted.] Dafny; the lemma's statement and the rest of its proof
  (only one assertion is missing); the deletion procedure's claim that the
  assertion is necessary (removing it breaks verification). The limitations
  call the setting optimistic: the scaffold and every other assertion of the
  same repository are available. \emph{Lines 892, 911.}
\item[Cheat surface.] The range does not say how a candidate is checked beyond
  acceptance (whether an added \dfy{assume}, a changed statement or an edited
  scaffold is refused). A full answer says so and names what one would need
  to know: the harness's checks on the edit. \emph{Lines 892, 911 (by
  absence).}
\item[Metric.] Solved if any of the first $k \le 10$ samples verifies; 56.6\%
  (82/145) with proof-similarity retrieval plus the placeholder, against 6.2\%
  for code only, 29.6\% with error feedback, 34.4\% with the placeholder alone
  and 51.7--55.0\% with other example selections; Cedar 16.6\% with the
  placeholder only against about 59\% with similar examples. \emph{Lines 890,
  892, 898--909.}
\item[Budget.] At most 10 samples per task, GPT-4o at temperature 1.0; the
  top 6 retrieved contexts become few-shot examples. \emph{Lines 892, 896.}
\item[Variance.] Not reported: one model, one run. \emph{Lines 892, 911.}
\item[Artifact.] DafnyGym (145 tasks over 33k lines of Dafny); the range does
  not say whether the benchmark or the tool is released. \emph{Line 892.}
\item[Ninth line (H3).] The verifier's error gives a \emph{where}, the
  postcondition or assertion that fails and its line, which Laurel's
  placeholder computes from the control-flow graph; it does not give the
  \emph{what}. In H3 the learner supplies the missing fact: an equality
  between two slices of the sequence (Dafny) or an extensional equality
  between the sequence and its split (Verus), which no message names. A
  full answer: where from the tool, what from the learner (or from a
  retrieved similar proof, Laurel's other half).
\end{description}

\end{document}
